% Especificacao de Caso de Uso: suc_gerenciar_usuarios
%
% Conversao do documento original do NPI (Google Docs, letter, Times New Roman
% no corpo e Arial Bold nos titulos). O layout segue o PDF antigo:
% cabecalho em tabela de 3 linhas, rodape "Confidencial", capa alinhada a
% direita e corpo recuado 1,2 cm em relacao a margem.

\documentclass[12pt,letterpaper]{article}

\usepackage[utf8]{inputenc}
\usepackage[T1]{fontenc}
\usepackage[brazilian]{babel}
\usepackage{mathptmx}                 % Times New Roman no corpo
\usepackage{helvet}                   % Arial nos titulos
\usepackage[letterpaper,left=2.54cm,right=2.54cm,top=4.1cm,bottom=2.6cm,
            headheight=58pt,headsep=12pt,footskip=30pt]{geometry}
\usepackage{graphicx}
\usepackage{float}
\usepackage{array}
\usepackage{fancyhdr}
\usepackage{titlesec}
\usepackage{tocloft}
\usepackage{enumitem}
\usepackage{changepage}

% ------------------------------------------------------------- metadados

% Os \allowbreak permitem quebrar o nome depois dos underscores. Sem eles a
% palavra e indivisivel e estoura a margem da linha que a contem.
\newcommand{\ucnome}{suc\_\allowbreak gerenciar\_\allowbreak usuarios}
\newcommand{\ucprojeto}{Reserva de salas}
\newcommand{\ucversao}{1.0}
\newcommand{\ucdata}{30/10/2026}
\newcommand{\ucautor}{José Ronaldo dos Santos Jr.}
\newcommand{\ucorgao}{Núcleo de Prática de Informática - NPI}
\newcommand{\uctitulo}{Especificação de Caso de Uso: \ucnome}

% O recuo do corpo em relacao a margem, medido no PDF original.
\newlength{\recuo}
\setlength{\recuo}{1.2cm}

\setlength{\parindent}{0pt}
\setlength{\parskip}{6pt}

% O nome do caso de uso e uma palavra indivisivel; sem folga extra ele estica
% os espacos da linha justificada que o contem.
\setlength{\emergencystretch}{3em}

% --------------------------------------------------- cabecalho e rodape

\pagestyle{fancy}
\fancyhf{}
\renewcommand{\headrulewidth}{0pt}
\renewcommand{\footrulewidth}{0pt}

\newcolumntype{H}[1]{>{\raggedright\arraybackslash}p{#1}}

\fancyhead[L]{%
  \hspace*{-\tabcolsep}%
  \begin{tabular}{|H{0.665\textwidth}|H{0.275\textwidth}|}
    \hline
    \textbf{\ucprojeto} & Versão:\qquad \ucversao \\
    \hline
    \uctitulo & Data:\ \ \ucdata \\
    \hline
    \uctitulo{} - \ucautor & \\
    \hline
  \end{tabular}%
}

\fancyfoot[L]{\small Confidencial}
\fancyfoot[C]{\small\parbox[t]{5.6cm}{\centering\copyright\ucorgao, 2026}}
\fancyfoot[R]{\small Página \thepage}

% ------------------------------------------------------ titulos de secao

% "1." nas secoes e "3.1" nas subsecoes, como no Word.
\renewcommand{\thesection}{\arabic{section}.}
\renewcommand{\thesubsection}{\arabic{section}.\arabic{subsection}}

% O numero fica na margem e o titulo alinha com o corpo recuado.
\titleformat{\section}[hang]
  {\sffamily\large\bfseries}
  {\makebox[\recuo][l]{\thesection}}{0pt}{}
\titleformat{\subsection}[hang]
  {\sffamily\normalsize\bfseries}
  {\makebox[\recuo][l]{\thesubsection}}{0pt}{}

\titlespacing*{\section}{0pt}{6pt}{2pt}
\titlespacing*{\subsection}{0pt}{4pt}{2pt}

% ------------------------------------------------------------- o indice

\addto\captionsbrazilian{\renewcommand{\contentsname}{Índice}}

\renewcommand{\cfttoctitlefont}{\hfill\sffamily\LARGE\bfseries}
\renewcommand{\cftaftertoctitle}{\hfill}
\setlength{\cftbeforetoctitleskip}{0pt}
\setlength{\cftaftertoctitleskip}{18pt}

% Entradas em Times, sem pontilhado, numeros nas mesmas colunas do corpo.
\renewcommand{\cftsecfont}{\rmfamily}
\renewcommand{\cftsecpagefont}{\rmfamily}
\renewcommand{\cftsubsecfont}{\rmfamily}
\renewcommand{\cftsubsecpagefont}{\rmfamily}
\renewcommand{\cftsecleader}{\hfill}
\renewcommand{\cftsubsecleader}{\hfill}
\renewcommand{\cftsecaftersnumb}{\hspace{0.3cm}}
\renewcommand{\cftsubsecaftersnumb}{\hspace{0.3cm}}
\setlength{\cftsecindent}{\recuo}
\setlength{\cftsecnumwidth}{1.1cm}
\setlength{\cftsubsecindent}{2.3cm}
\setlength{\cftsubsecnumwidth}{1.1cm}
\setlength{\cftbeforesecskip}{14pt}
\setlength{\cftbeforesubsecskip}{4pt}

% --------------------------------------------- corpo recuado e os fluxos

% Paragrafos comuns do corpo entram recuados; os titulos de secao usam
% \hspace negativo embutido no \titleformat para voltar a margem.
\newcommand{\corpo}{\setlength{\leftskip}{\recuo}}

% O rotulo "1." cai na mesma coluna do corpo recuado (listas ignoram
% \leftskip, por isso a margem entra explicita aqui).
\newlist{fluxo}{enumerate}{1}
\setlist[fluxo]{label=\arabic*., align=left, labelwidth=1.2em, labelsep=0.4em,
                leftmargin=\dimexpr\recuo+1.6em\relax, labelindent=0pt,
                topsep=6pt, itemsep=6pt, parsep=0pt}

\begin{document}

% ------------------------------------------------------------------ capa

\thispagestyle{empty}
\vspace*{-2.2cm}
{\sffamily\bfseries
  \hrule height 0.8pt
  \vspace{10pt}
  \raggedleft \Large \ucorgao \par
  \vspace{10pt}
  \hrule height 0.8pt
  \vspace{30pt}
  \raggedleft \Large \ucprojeto \par
  \vspace{10pt}
  \raggedleft \Large \uctitulo \par
  \vspace{34pt}
  \raggedleft \large Versão \ucversao \par
}
\newpage

% -------------------------------------------- historico da revisao e indice

\begin{center}
  {\sffamily\LARGE\bfseries Histórico da Revisão}
\end{center}

\vspace{6pt}

\hspace*{-\tabcolsep}%
\begin{tabular}{|p{0.222\textwidth}|p{0.102\textwidth}|p{0.352\textwidth}|p{0.222\textwidth}|}
  \hline
  \centering\textbf{Data} & \centering\textbf{Versão} & \centering\textbf{Descrição}
    & \centering\arraybackslash\textbf{Autor} \\
  \hline
  \ucdata & \ucversao & Primeira versão & \ucautor \\
  \hline
  & & & \\
  \hline
  & & & \\
  \hline
  & & & \\
  \hline
\end{tabular}

\vspace{18pt}

\begin{adjustwidth}{0pt}{1.2cm}
  \tableofcontents
\end{adjustwidth}
% \tableofcontents forca \thispagestyle{plain}; devolvemos cabecalho e rodape.
\thispagestyle{fancy}
\newpage

% ------------------------------------------------------------- o documento

\begin{center}
  {\sffamily\Large\bfseries \uctitulo}
\end{center}

\corpo

O Objetivo deste documento é apresentar todas as especificações dos requisitos
funcionais do caso de uso \ucnome.

\section{Breve Descrição}

A função deste caso de uso é permitir que o usuário do sistema, coordenador,
seja capaz de visualizar, criar, editar, e excluir usuários.

\section{Fluxo Básico de Eventos}

\begin{fluxo}
  \item O usuário acessa o menu principal.
  \item O usuário acessa a seção de gerenciamento de usuários
  \item O usuário escolhe a opção de cadastrar usuário.
  \item O sistema exibe um formulário de cadastro de usuário.
  \item O usuário faz a inserção dos dados.
  \item O usuário salva o cadastro de um novo usuário.
  \item O sistema valida e salva os dados recebidos e finaliza o processo de
        cadastro.
\end{fluxo}

\section{Fluxos Alternativos}

Nesta seção, serão abordados fluxos alternativos relacionados ao caso de uso
\ucnome, um fluxo alternativo representa um comportamento alternativo
geralmente derivados de exceções que ocorrem no fluxo principal.

\subsection{Visualizar Usuários}

\begin{fluxo}
  \item O usuário acessa o menu principal.
  \item O usuário acessa a seção de gerenciamento de usuários.
  \item O usuário seleciona um usuário.
  \item O usuário visualiza as informações desejadas.
\end{fluxo}

\subsection{Editar Usuários}

\begin{fluxo}
  \item O usuário acessa o menu principal.
  \item O usuário acessa a seção de gerenciamento de usuários.
  \item O usuário faz a alteração dos dados.
  \item O sistema valida e salva os dados recebidos e finaliza o processo.
\end{fluxo}

\subsection{Excluir Usuários}

\begin{fluxo}
  \item O usuário acessa o menu principal.
  \item O usuário acessa a seção de gerenciamento de usuários.
  \item O usuário seleciona o usuário que deseja excluir.
  \item O usuário seleciona a opção de excluir.
  \item Os dados são excluídos do sistema.
\end{fluxo}

\subsection{Tratamento de Dados}

O sistema valida os dados inseridos pelo usuário. O sistema apresenta uma
mensagem de erro caso os dados inseridos pelo usuário não estejam adequados
para o sistema.

\section{Subfluxos}

Não existem subfluxos associados a este caso de uso.

\section{Cenários Chave}

Os cenários chaves deste caso de uso são: criação, visualização, edição e
exclusão de usuários cadastrados no sistema.

\section{Condições Prévias}

Não se aplica.

\section{Condições Posteriores}

Nesta seção serão listados os possíveis estados do sistema imediatamente após
a conclusão do caso de uso em questão.

\subsection{Cadastrar Usuário}

Um usuário será cadastrado no sistema e salvo no banco de dados.

\subsection{Editar Usuário}

Um usuário existente terá suas informações atualizadas com novas informações
no banco de dados.

\subsection{Deletar Usuário}

Um usuário existente será apagado do banco de dados.

\section{Pontos de Extensão}

O caso de uso \ucnome{} não possui pontos de extensão.

\section{Requisitos Especiais}

O caso de uso \ucnome{} não possui requisitos especiais.

\section{Informações Adicionais}

Na Figura 1 são apresentados os atores e funcionalidades do caso de uso.

\begin{figure}[H]
  \centering
  Figura 1 -- Caso de uso

  \vspace{10pt}
  \includegraphics[width=0.6\textwidth]{../../../out/src/plantuml/casos-de-uso/suc_gerenciar_usuarios.png}

  \vspace{10pt}
  Fonte: O autor.
\end{figure}

\end{document}
